\section{系统整合}

\subsection{文件夹结构}

作者建议将所有 Cowork 上下文文件组织在一个统一的文件夹中，并将 Cowork 指向该文件夹。这样每次对话都能自动获取完整的上下文。

合理的文件夹结构确保：
\begin{itemize}
    \item 所有上下文文件集中管理，便于维护
    \item Cowork 能一次性读取全部上下文，而非逐个指定
    \item 新增上下文文件时只需放入文件夹即可生效
\end{itemize}

\subsection{Meta-Prompt --- 防错启动提示}

在运行任何工作流之前，先用一个 meta-prompt 对 Cowork 进行预设。作者声称这个单一提示能预防 90\% 用户常抱怨的错误。

\begin{warningbox}{常见陷阱}
很多用户跳过 meta-prompt 直接运行工作流，导致 Cowork 缺乏必要的约束和偏好信息，输出质量不稳定。Meta-prompt 的作用类似于 system prompt，为后续所有交互设定基调和规则。
\end{warningbox}

\subsection{本章小结}

系统整合的两个关键：（1）统一的文件夹结构确保上下文完整性；（2）meta-prompt 设定全局规则，预防常见错误。两者结合让 Cowork 从零散的工具集合变成一个有机的工作系统。

